\documentclass[12pt,a4paper]{article}

\usepackage[top=2.5cm,bottom=2.5cm,left=2.5cm,right=2.5cm]{geometry}
\usepackage{graphicx}
\usepackage{listings}
\usepackage{xcolor}
\usepackage{hyperref}
\usepackage{fancyhdr}
\usepackage{booktabs}
\usepackage{array}
\usepackage{enumitem}
\usepackage{parskip}
\usepackage{float}
\usepackage{caption}
\usepackage{subcaption}
\usepackage{titlesec}
\usepackage{microtype}
\usepackage{setspace}
\usepackage{tcolorbox}
\tcbuselibrary{skins}

\definecolor{primary}{RGB}{25,70,140}
\definecolor{accent}{RGB}{200,150,20}
\definecolor{codebg}{RGB}{246,248,250}
\definecolor{codeframe}{RGB}{180,200,230}
\definecolor{kw}{RGB}{25,70,140}
\definecolor{str}{RGB}{150,40,150}
\definecolor{cmt}{RGB}{80,140,80}
\definecolor{gray}{RGB}{120,120,120}

\lstdefinestyle{code}{
  backgroundcolor=\color{codebg},
  basicstyle=\ttfamily\small,
  keywordstyle=\color{kw}\bfseries,
  stringstyle=\color{str},
  commentstyle=\color{cmt}\itshape,
  numberstyle=\tiny\color{gray},
  numbers=left, numbersep=6pt,
  breaklines=true, frame=single,
  rulecolor=\color{codeframe},
  xleftmargin=12pt, tabsize=4,
  showstringspaces=false,
  captionpos=b,
}
\lstset{style=code}

\titleformat{\section}{\Large\bfseries\color{primary}}{\thesection}{0.8em}{}[\vspace{-4pt}\color{primary}\hrule\vspace{4pt}]
\titleformat{\subsection}{\large\bfseries\color{primary!80!black}}{\thesubsection}{0.7em}{}
\titleformat{\subsubsection}{\normalsize\bfseries\color{primary!60!black}}{\thesubsubsection}{0.6em}{}

\pagestyle{fancy}\fancyhf{}
\fancyhead[L]{\small\color{primary}\textbf{Mini Game Hub}}
\fancyhead[R]{\small\color{gray}SSL Project Report}
\fancyfoot[C]{\small\thepage}
\renewcommand{\headrulewidth}{0.4pt}

\newcolumntype{L}[1]{>{\raggedright\arraybackslash}p{#1}}
\newcolumntype{C}[1]{>{\centering\arraybackslash}p{#1}}

\newtcolorbox{notebox}{
  colback=codebg, colframe=primary!40,
  boxrule=0.5pt, arc=4pt, left=6pt, right=6pt, top=4pt, bottom=4pt,
}

\hypersetup{colorlinks,linkcolor=primary,urlcolor=primary}

\begin{document}

\begin{titlepage}
\centering\vspace*{2.5cm}
{\color{primary}\rule{\linewidth}{1.5pt}}\vspace{6pt}
{\color{primary}\rule{\linewidth}{0.4pt}}\vspace{1.2cm}

{\Huge\bfseries\color{primary} Mini Game Hub}\\\vspace{0.4cm}
{\large\itshape SSL Project Technical Report}\\\vspace{1cm}

{\color{primary}\rule{\linewidth}{0.4pt}}\vspace{6pt}
{\color{primary}\rule{\linewidth}{1.5pt}}\vspace{2cm}

\begin{tcolorbox}[colback=codebg,colframe=primary!50,arc=5pt,width=0.82\textwidth]
\centering\large
A secure, multi-user game hub featuring four classic\\
two-player board games, built with Bash, Python,\\
Pygame, and NumPy.
\end{tcolorbox}

\vspace{2cm}
{\large\textbf{Team:} Rajit \textbullet\ Aagam}\\[0.5cm]
{\large\textbf{Date:} April 2026}\\[0.5cm]
{\large\textbf{Stack:} Python 3 \textbullet\ Pygame \textbullet\ NumPy \textbullet\ Bash}
\vfill
{\color{gray}\small Systems and Software Lab (SSL) Academic Project}
\end{titlepage}

\tableofcontents
\newpage

\section{External Libraries}

\begin{table}[H]
\centering
\caption{External libraries and tools used}
\begin{tabular}{L{3cm} L{3cm} L{8cm}}
\toprule
\textbf{Library / Tool} & \textbf{Version} & \textbf{Purpose} \\
\midrule
Python 3  & 3.10+  & Primary application language \\
Pygame    & 2.x    & Window management, event loop, sprite/image rendering, font drawing \\
NumPy     & 1.x    & 2-D array board representation and vectorised win-detection \\
Bash      & 5.x    & Authentication front-end: password hashing, TSV credential lookup, app launch \\
\texttt{csv} (stdlib) & — & Reading and writing \texttt{history.csv} game records \\
\texttt{datetime} (stdlib) & — & Timestamping match results \\
\texttt{sys} (stdlib) & — & Access to interpreter-level functionality via command-line arguments\\
\texttt{os} (stdlib) & — & Operating system interface: file/directory management, path handling, and process control\\
\bottomrule
\end{tabular}
\end{table}

No third-party packages beyond Pygame and NumPy are required.

% ════════════════════════════════════════════════════════════════
\section{Design \& Architecture}
% ════════════════════════════════════════════════════════════════

\subsection{High-Level Structure}

The system is divided into three layers with clear responsibilities:

\begin{enumerate}[leftmargin=2em]
  \item \textbf{Authentication layer} (Bash): Runs in the terminal before any graphical window opens. Handles user registration and login using SHA-256 password hashing and a plain-text TSV credential store, then launches the Python application with the two player names as arguments.
  \item \textbf{Hub / controller layer} (\texttt{gaming.py}): Owns the single shared Pygame surface (1000$\times$700 px). Manages the main menu, game-selection screen, post-game flow, leaderboard, and CSV persistence. All UI transitions are handled here.
  \item \textbf{Game layer} (individual game modules): Each game is a self-contained class that receives the shared surface and a set of controller-class references via its constructor. The class runs its own \texttt{run()} event loop and calls back into the hub classes on termination.
\end{enumerate}

\subsection{Key Design Decisions}

\textbf{Single shared Pygame surface.}
Rather than creating a new window per game, all classes draw to the same 1000$\times$700 surface. This avoids multi-window complexity and gives a seamless transition feel.

\textbf{Constructor injection instead of globals.}
Every game class receives \texttt{GameSelected}, \texttt{Resign}, \texttt{CommonWC}, and \texttt{UpdateCSV} as constructor parameters. This keeps each module independently testable and decoupled from \texttt{gaming.py}'s internal state.

\textbf{Flat file persistence.}
A CSV for match history and a TSV for credentials were chosen over a database because they require no external server, are human-readable, and are trivially portable. Python's built-in \texttt{csv} module handles all I/O.

\textbf{Move array pattern.}
Each game appends every move as a tuple to a shared \texttt{movearray} list passed in at construction. The first element is always a sentinel \texttt{(0, "GameName", 0)}. On game end, the entire array is optionally serialised to \texttt{SavedGames.txt}.

% ════════════════════════════════════════════════════════════════
\section{Features Implemented}
% ════════════════════════════════════════════════════════════════

\subsection{Authentication System}
\begin{itemize}[leftmargin=2em]
  \item User registration and login via terminal (Bash).
  \item Passwords hashed with SHA-256; never stored in plain text.
  \item Credentials persisted in \texttt{users.tsv} (tab-separated: username, hash).
  \item Input masking (\texttt{read -s}) during password entry.
\end{itemize}

\subsection{Main Menu}
\begin{itemize}[leftmargin=2em]
  \item Custom background image with semi-transparent player name overlays.
  \item Navigation buttons for Leaderboard, Stats, How To Play, Settings, and four game tiles.
  \item Game tiles rendered as clickable icon regions; hovering and clicking routed to \texttt{GameSelected}.
\end{itemize}

\subsection{Games}

\subsubsection{Tic-Tac-Toe}
Modified version of the classic game, this one is played on a $10\times10$ board. Board stored as a NumPy array; win detection checks row/column/diagonal sums. Draws detected when all cells are filled with no winner.

\subsubsection{Connect Four}
7$\times$7 dropping-disc game. Gravity implemented by scanning from the bottom of each column for the first empty cell. Win detection slides a window of four cells across all rows, columns, and both diagonals.

\subsubsection{Othello}
8$\times$8 disc-flipping strategy game. Valid move generation traces all eight directions from a candidate cell. On placement, all flippable opponent discs in every direction are converted. Auto-passes turn when no moves exist. Disc count determines winner.

\subsubsection{Catan (Simplified)}
The most complex game: a two-player adaptation of \textit{Settlers of Catan} on a 7-tile hexagonal board.

\begin{itemize}[leftmargin=2em]
  \item \textbf{Hex board generation:} Axial coordinate system with random resource and number-token placement. Port system assigned to ring tiles (2:1 specific-resource and 3:1 general ports).
  \item \textbf{Graph construction:} Vertices and edges derived from hex geometry; shared vertices between adjacent tiles deduplicated via rounded grid keys. Adjacency map supports road/settlement validation.
  \item \textbf{Setup phase:} Alternating settlement+road placement with distance-rule enforcement.
  \item \textbf{Main phase:} Dice roll, resource distribution to adjacent settlements/cities, building (road, settlement, city), development card play, and end-turn.
  \item \textbf{Robber:} Triggered on roll of 7; active player moves robber, players with 7+ cards discard half, player steals one resource from opponent on that tile.
  \item \textbf{Development cards:} Knight, Road Building, Monopoly, Year of Plenty.
  \item \textbf{Trading:} Maritime bank trade (4:1 standard), Port trades (2:1/3:1) and player-to-player trade with accept/reject flow.
  \item \textbf{Longest Road / Largest Army:} Computed via DFS over the road graph; 2 VP bonus awarded dynamically.
  \item \textbf{Win condition:} First to 6 Victory Points (settlements, cities, special cards).
\end{itemize}

\subsection{Post-Game Flow}
\begin{itemize}[leftmargin=2em]
  \item \textbf{Win screen}: \texttt{WinnerFace.png} with the winner's name rendered in large type.
  \item \textbf{Draw screen}: dedicated \texttt{ItsaDraw.png} layout.
  \item \textbf{Resignation}: Confirmation screen (\texttt{ResignCat.png}); move array updated before result is logged.
  \item Four post-game actions: \textit{Save game} (serialise move array to \texttt{SavedGames.txt}), \textit{Rematch} (reinvoke same game/mode), \textit{Stats} (show plots of player and game stats), \textit{Return} (back to \texttt{GameSelected}).
\end{itemize}

\subsection{Leaderboard}
\begin{itemize}[leftmargin=2em]
  \item Reads \texttt{history.csv} and aggregates wins, losses, and win/loss ratio per player.
  \item Draws award 0.5 wins and 0.5 losses to each player.
  \item Filterable by game (Overall, Tic-Tac-Toe, Connect Four, Othello).
  \item Sortable by wins, losses, or ratio via in-screen dropdown.
  \item Infinite ratio (zero losses) displayed as $\infty$.
\end{itemize}

\subsection{How To Play}
In-hub rules reference accessible from the main menu. Three game thumbnails navigate to per-game instruction pages rendered on a clean background.

\subsection{Persistent Data}
\begin{itemize}[leftmargin=2em]
  \item Every game result written to \texttt{history.csv} with serial number, timestamp, game name, both players, and result code.
  \item Serial counter auto-incremented in \texttt{Serial.txt}.
  \item Move arrays optionally saved to \texttt{SavedGames.txt}.
\end{itemize}

% ════════════════════════════════════════════════════════════════
\section{Directory Structure}
% ════════════════════════════════════════════════════════════════

\begin{lstlisting}[language=bash, caption={Annotated project tree}]
SSL-Project/
|
|-- main.sh                # Entry point: Bash auth, launches gaming.py
|-- gaming.py              # Hub controller: menus, leaderboard, CSV, UI classes
|-- Catan.py               # Full Catan game implementation (1057 lines)
|-- Redirect.py            # Alternate lightweight menu/redirect module
|
|-- users.tsv              # Credential store: username<TAB>sha256-hash
|-- history.csv            # Match log: serial,timestamp,game,p1,p2,result
|-- SavedGames.txt         # Serialised move arrays from completed games
|-- Serial.txt             # Auto-incrementing match ID counter
|
|-- MCA Background.png     # Main menu background
|-- Connect4Background.png # Connect Four game background
|-- OthelloBackground.png  # Othello game background
|-- WinnerFace.png         # Win result screen background
|-- ResignCat.png          # Resignation confirmation background
|-- MagicCatFace.png       # Draw/misc result background
|
|-- settlement_p1.png      # Catan: Player 1 settlement sprite
|-- settlement_p2.png      # Catan: Player 2 settlement sprite
|-- road_p1.png            # Catan: Player 1 road sprite
|-- road_p2.png            # Catan: Player 2 road sprite
|-- wood.png               # Catan: Wood resource tile
|-- brick.png              # Catan: Brick resource tile
|-- wool.png               # Catan: Wool resource tile
|-- wheat.png              # Catan: Wheat resource tile
|-- ore.png                # Catan: Ore resource tile
\end{lstlisting}

% ════════════════════════════════════════════════════════════════
\section{Installation \& Setup}
% ════════════════════════════════════════════════════════════════

\subsection{Prerequisites}
\begin{itemize}[leftmargin=2em]
  \item Python 3.10 or later
  \item pip
  \item A Unix-like shell (Linux or macOS); on Windows, use WSL
\end{itemize}

\subsection{Installing Dependencies}

\begin{lstlisting}[language=bash, caption={Install Python dependencies}]
pip install -r requirements.txt
\end{lstlisting}

\noindent The \texttt{requirements.txt} file contains:

\begin{lstlisting}[caption={requirements.txt}]
pygame>=2.0.0
numpy>=1.21.0
\end{lstlisting}

\subsection{Font Asset}
The hub uses \texttt{Fredoka\_Expanded-Bold.ttf}. Place this font file in the project root directory (same level as \texttt{gaming.py}). It can be downloaded from Google Fonts:

\begin{notebox}
\textbf{Font:} \url{https://fonts.google.com/specimen/Fredoka}
\end{notebox}

\subsection{Large Assets}
All PNG background and sprite assets are bundled in the submission zip. No external downloads are needed for images.

\subsection{Using the Makefile}

\begin{lstlisting}[language=bash, caption={Makefile targets}]
make install   # pip install -r requirements.txt
make run       # bash main.sh
make clean     # remove __pycache__ and .pyc files
\end{lstlisting}

% ════════════════════════════════════════════════════════════════
\section{Usage}
% ════════════════════════════════════════════════════════════════

\subsection{Running the Hub}

\begin{lstlisting}[language=bash, caption={Launching the game hub}]
bash main.sh
\end{lstlisting}

The terminal will prompt for two players to log in (or register). On successful authentication, the Pygame window opens automatically.

\subsection{Navigation}

\begin{enumerate}[leftmargin=2em]
  \item \textbf{Main Menu}: Click a game icon at the bottom to select a game. Use the top navigation bar for Leaderboard, Stats, How To Play, and Settings.
  \item \textbf{In-Game}: Click to place pieces / make moves. A Resign button is available at any time.
  \item \textbf{Post-Game}: Choose Save, Rematch, Stats or Return from the result screen.
  \item \textbf{Leaderboard}: Select Overall or a specific game; use the Sort dropdown to reorder by wins, losses, or ratio.
\end{enumerate}


% ════════════════════════════════════════════════════════════════
\section{Retrospective}
% ════════════════════════════════════════════════════════════════

\subsection{Challenges and How We Overcame Them}

\textbf{Hexagonal grid mathematics (Catan).}
Computing vertex positions from axial coordinates and correctly deduplicating shared vertices between adjacent hexagons required careful floating-point rounding. We resolved this by normalising each vertex coordinate to a rounded grid key \texttt{(round(x/10), round(y/10))} before storing, guaranteeing that geometrically identical points map to the same ID regardless of which tile generated them.

\textbf{Graph construction for road validation.}
Validating road placement (must connect to existing network) and computing the longest road required a proper adjacency graph. We built the edge-to-vertex map dynamically during board initialisation and used iterative DFS with backtracking to compute road length, which also handles branching paths correctly.

\textbf{State transitions across UI classes.}
Each class runs its own Pygame event loop, making it easy to lose track of which loop is ``active''. We standardised the pattern: every class returns cleanly from \texttt{run()} when it wants to hand control back, and all navigation is done by the \textit{caller} after \texttt{run()} returns - no global state, no threading.

\textbf{CSV header inconsistency.}
Early in development some rows were written without headers, causing \texttt{DictReader} field misalignment. We fixed this by writing the header once during file creation and using \texttt{newline=""} consistently to avoid platform-specific line-ending issues.

\textbf{Player-to-player trading in Catan.}
Implementing the ``awaiting response'' pattern, pausing one player's input while the other accepts or rejects a trade, required an explicit \texttt{awaiting\_response} flag that gates all normal input processing. This proved simpler than any callback or threading approach.

\subsection{What We Would Do Differently}

\begin{itemize}[leftmargin=2em]
  \item \textbf{Add password salting.} The current SHA-256 scheme is unsalted, meaning identical passwords produce identical hashes. A per-user random salt stored alongside the hash would eliminate rainbow-table vulnerability.
  \item \textbf{Implement Chess.} The game map reserves slot 3 for Chess but the class was never built. Given more time, integrating \texttt{python-chess} for move validation and a minimax AI would complete the original five-game vision.
  \item \textbf{Automated testing.} Because every class is tightly coupled to Pygame's display, unit-testing game logic required running the GUI. Separating board/logic classes from rendering classes from the start would have enabled clean pytest coverage.
  \item \textbf{Config file for assets.} Font paths and image paths are hardcoded strings scattered across multiple files. A single \texttt{config.py} with all asset paths would make the project far easier to port or rename assets.
\end{itemize}

% ════════════════════════════════════════════════════════════════
\section*{References}
\addcontentsline{toc}{section}{References}
% ════════════════════════════════════════════════════════════════

\begin{enumerate}[leftmargin=2em]
  \item Pygame documentation: \url{https://www.pygame.org/docs/}
  \item NumPy documentation: \url{https://numpy.org/doc/}
  \item \textit{Settlers of Catan} rulebook: Catan GmbH, used as reference for game mechanics.
  \item Google Fonts - Fredoka typeface: \url{https://fonts.google.com/specimen/Fredoka}
  \item SHA-256 in Bash: \texttt{man sha256sum} / \texttt{man openssl dgst}
\end{enumerate}

\end{document}
